\documentclass[../../main.tex]{subfiles}
\begin{document}

{\bf [解答]}

% TODO: 概要図

題意の曲線$y=f(x)$とすると，倍角公式，三倍角の公式より
\begin{align}
  f'(x)&=6\cos 2x-3\sin 3x  \nonumber\\
  &=6(1-2\sin^2 x)-3(3\sin x-4\sin^3x)  \nonumber\\
  &=3(4\sin^3x-4\sin^2x-3\sin x+2) \label{eq:1}
\end{align}
である．題意の直線の傾きは$-3$だから，$f'(x)=-3$となる$x$での接線を求めればよい．
\cref{eq:1}から
\begin{align}
  & f'(x)=-3
  \Longleftrightarrow
  & 4\sin^3x-4\sin^2x-3\sin x+2=-1 \\
  \Longleftrightarrow
  & (4\sin^2x-3)(\sin x-1)=0  \\
  \Longleftrightarrow
  &\sin x=\dfrac{\sqrt{3}}{2},1 \quad(\because 0<x<\pi)\\
  \Longleftrightarrow
  & x=\dfrac{\pi}{2},\dfrac{\pi}{3},\dfrac{2\pi}{3}
\end{align}
である．
それぞれ$f(x)$に代入すると$f(\pi/2)=0$，$f(\pi/3)=3\sqrt{3}/2-1$，$f(2\pi/3)=1-3\sqrt{3}/2$だから求める方程式は
\begin{align}
  \begin{cases}
    y=-3(x-\dfrac{\pi}{2})  \\
    y=-3(x-\dfrac{\pi}{3})+\dfrac{3\sqrt{3}}{2}-1  \\
    y=-3(x-\dfrac{2\pi}{3})+1-\dfrac{3\sqrt{3}}{2}
  \end{cases} \tag{答}
\end{align}
となる．
\newpage
\end{document}
